\documentclass[12pt,a4paper]{article}

\usepackage[utf8]{inputenc}
\usepackage[vietnamese]{babel}
\usepackage{amsmath,amssymb,amsthm}
\usepackage{geometry}
\usepackage{enumitem}
\usepackage[most]{tcolorbox}
\usepackage{hyperref}
\usepackage{fancyhdr}
\usepackage{graphicx}
\usepackage{tikz}
\usepackage{eso-pic}
\usepackage{xcolor}

\geometry{a4paper, margin=2cm, bottom=2.5cm}
\hypersetup{colorlinks=true, linkcolor=blue!70!black, urlcolor=blue}
\setlist[itemize]{topsep=4pt, itemsep=2pt}
\setlist[enumerate]{topsep=4pt, itemsep=2pt}

\AddToShipoutPictureFG{%
  \AtPageCenter{%
    \begin{tikzpicture}[remember picture, overlay]
      \node[rotate=45, scale=1, opacity=0.06]{\fontsize{60}{72}\selectfont\bfseries\sffamily Đặng Tấn Quí};
    \end{tikzpicture}%
  }%
}

\pagestyle{fancy}
\fancyhf{}
\fancyhead[L]{\small\textit{Đại số sơ cấp}}
\fancyhead[R]{\small\textit{Đặng Tấn Quí}}
\fancyfoot[C]{\thepage}
\fancyfoot[L]{\footnotesize\textcopyright\ Đặng Tấn Quí}
\fancyfoot[R]{\footnotesize SĐT: 0377\,122\,210}
\renewcommand{\headrulewidth}{0.4pt}
\renewcommand{\footrulewidth}{0.4pt}

\fancypagestyle{plain}{\fancyhf{}\fancyfoot[C]{\thepage}\fancyfoot[L]{\footnotesize\textcopyright\ Đặng Tấn Quí}\fancyfoot[R]{\footnotesize SĐT: 0377\,122\,210}\renewcommand{\headrulewidth}{0pt}\renewcommand{\footrulewidth}{0.4pt}}

\newtcolorbox{dinhnghia}[1][]{enhanced, breakable, colback=blue!3!white, colframe=blue!60!black, fonttitle=\bfseries, title={Định nghĩa}, #1}
\newtcolorbox{dinhly}[1][]{enhanced, breakable, colback=green!3!white, colframe=green!50!black, fonttitle=\bfseries, title={Định lý}, #1}
\newtcolorbox{tinhchat}[1][]{enhanced, breakable, colback=yellow!5!white, colframe=orange!50!black, fonttitle=\bfseries, title={Tính chất}, #1}
\newtcolorbox{phuongphap}[1][]{enhanced, breakable, colback=gray!5!white, colframe=gray!50!black, fonttitle=\bfseries, title={Phương pháp}, #1}
\newtcolorbox{luuy}[1][]{enhanced, breakable, colback=red!3!white, colframe=red!50!black, fonttitle=\bfseries, title={Lưu ý}, #1}
\newtcolorbox{congthuc}[1][]{enhanced, breakable, colback=cyan!3!white, colframe=cyan!50!black, fonttitle=\bfseries, title={Công thức}, #1}
\newtcolorbox{vidu}[1][]{enhanced, breakable, colback=violet!3!white, colframe=violet!50!black, fonttitle=\bfseries, title={Ví dụ}, #1}

\begin{document}

\thispagestyle{empty}
\begin{tikzpicture}[remember picture, overlay]
  \fill[blue!80!black] (current page.south west) rectangle (current page.north east);
  \node[anchor=west, white, font=\fontsize{26}{32}\bfseries\selectfont]
    at ([xshift=3cm, yshift=-7.5cm]current page.north west) {ĐẠI SỐ SƠ CẤP};
  \node[anchor=west, white, font=\fontsize{18}{24}\selectfont]
    at ([xshift=3cm, yshift=-9cm]current page.north west) {Tập hợp, ánh xạ, số phức, cấu trúc đại số --- Năm 1, HK1};
  \node[anchor=west, white, font=\Large\bfseries] at ([xshift=3cm, yshift=-13cm]current page.north west) {Biên soạn:};
  \node[anchor=west, yellow!80!white, font=\fontsize{22}{28}\bfseries\selectfont] at ([xshift=3cm, yshift=-14.5cm]current page.north west) {ĐẶNG TẤN QUÍ};
  \node[anchor=south east, white, font=\Large\bfseries] at ([xshift=-2cm, yshift=3cm]current page.south east) {\the\year};
\end{tikzpicture}

\newpage
\tableofcontents
\newpage

%% ================================================================
%%  CHƯƠNG 1: QUY NẠP TOÁN HỌC
%% ================================================================
\section{Quy nạp toán học}

\begin{phuongphap}[title={Nguyên lý quy nạp}]
  Để chứng minh mệnh đề $P(n)$ đúng với mọi $n \in \mathbb{N}$, $n \ge n_0$:
  \begin{enumerate}
    \item \textbf{Bước cơ sở}: Chứng minh $P(n_0)$ đúng.
    \item \textbf{Bước quy nạp}: Chứng minh $\forall k \ge n_0$: $P(k) \Rightarrow P(k+1)$.
  \end{enumerate}
\end{phuongphap}

\begin{vidu}
  Chứng minh $1 + 2 + \cdots + n = \dfrac{n(n+1)}{2}$:
  \begin{itemize}
    \item Cơ sở: $n=1$: $1 = \dfrac{1 \cdot 2}{2}$ đúng.
    \item Quy nạp: Giả sử đúng với $n=k$. Khi đó $1 + \cdots + k + (k+1) = \dfrac{k(k+1)}{2} + (k+1) = \dfrac{(k+1)(k+2)}{2}$ đúng với $n=k+1$.
  \end{itemize}
\end{vidu}

\begin{vidu}
  Chứng minh $1^2 + 2^2 + \cdots + n^2 = \dfrac{n(n+1)(2n+1)}{6}$. Bất đẳng thức Bernoulli: $(1+x)^n \ge 1 + nx$ với $x \ge -1$, $n \in \mathbb{N}$.
\end{vidu}

%% ================================================================
%%  CHƯƠNG 2: SỐ PHỨC
%% ================================================================
\section{Số phức}

\subsection{Dạng đại số}

\begin{dinhnghia}
  \textbf{Số phức} (dạng chính tắc) $z = a + bi$, với $a, b \in \mathbb{R}$, $i^2 = -1$.
  \begin{itemize}
    \item $a = \mathrm{Re}(z)$: phần thực.
    \item $b = \mathrm{Im}(z)$: phần ảo.
    \item Tập các số phức ký hiệu $\mathbb{C}$.
    \item Tập số phức $\mathbb{C}$ có 2 luật thành trong: cộng ($+$) và nhân ($\cdot$)
  \end{itemize}
\end{dinhnghia}

\begin{congthuc}[title={Các phép toán}]
  $(a+bi) \pm (c+di) = (a \pm c) + (b \pm d)i$.
  
  $(a+bi)(c+di) = (ac-bd) + (ad+bc)i$.
  
  $\dfrac{a+bi}{c+di} = \dfrac{(a+bi)(c-di)}{c^2+d^2}$.

  $a + bi = a' + bi \Leftrightarrow a = a' \text{ và } b = b'$
\end{congthuc}

\begin{vidu}
  $(1+2i)(3-i) = 3 - i + 6i - 2i^2 = 5 + 5i$. $\dfrac{1}{i} = \dfrac{1 \cdot (-i)}{i \cdot (-i)} = \dfrac{-i}{1} = -i$. $|3+4i| = \sqrt{9+16} = 5$.
\end{vidu}

\begin{tinhchat}[title={$\mathbb{C}$ là trường}]
  \begin{itemize}
    \item Phần tử trung hòa của luật cộng: $(0, 0)$.
    \item Số phức đối của $(a, b)$: $(-a, -b) = -(a, b)$.
    \item Phần tử trung hòa của luật nhân: $(1, 0)$.
    \item Số phức nghịch đảo của $(a, b) \ne (0, 0)$: $\left(\dfrac{a}{a^2 + b^2}, \dfrac{-b}{a^2 + b^2}\right) = (a, b)^{-1}$.
  \end{itemize}
\end{tinhchat}

\begin{dinhnghia}
  \textbf{Mặt phẳng phức}: Mặt phẳng $Oxy$ với điểm $M = (a, b)$. Số phức $z = a + bi$ gọi là \textbf{tọa vị} của $M$. Trục $Ox$ (số phức $(a, 0)$): \textbf{trục thực}. Trục $Oy$ (số phức $(0, b)$): \textbf{trục ảo}. Số thực là trường hợp riêng của số phức. \textbf{Số ảo thuần túy} $(0, b)$ ($b \ne 0$) có bình phương luôn âm. \textbf{Đơn vị ảo} $i = (0, 1)$.
\end{dinhnghia}

\subsection{Dạng lượng giác}

\begin{dinhnghia}
  $z = a + bi = r(\cos\varphi + i\sin\varphi)$, với $r = |z| = \sqrt{a^2 + b^2}$, $\varphi = \arg z$ (argument).
\end{dinhnghia}

\begin{tinhchat}
  $r(\cos\varphi + i\sin\varphi) = r'(\cos\varphi' + i\sin\varphi') \Leftrightarrow r = r' \text{ và } \varphi = \varphi' + k2\pi$.
\end{tinhchat}

\begin{congthuc}[title={Nhân và chia}]
  $z_1 z_2 = r_1 r_2\bigl(\cos(\varphi_1 + \varphi_2) + i\sin(\varphi_1 + \varphi_2)\bigr)$.
  
  $\dfrac{z_1}{z_2} = \dfrac{r_1}{r_2}\bigl(\cos(\varphi_1 - \varphi_2) + i\sin(\varphi_1 - \varphi_2)\bigr)$.
\end{congthuc}

\begin{congthuc}[title={Công thức Moivre}]
  $z^n = r^n(\cos n\varphi + i\sin n\varphi)$ với $n \in \mathbb{Z}$.
\end{congthuc}

\begin{congthuc}[title={Tổng và bất đẳng thức tam giác}]
  $z_1 + z_2 = (a + a') + (b + b')i$. Điểm $A(a, b)$, $B(a', b')$, $C(a + a', b + b')$: vectơ $\overrightarrow{OC} = \overrightarrow{OA} + \overrightarrow{OB}$ (quy tắc hình bình hành).
  \[
    |z_1 + z_2| \le |z_1| + |z_2| \quad \text{(bất đẳng thức tam giác)}.
  \]
\end{congthuc}

\begin{dinhnghia}
  Số phức liên hợp: $\overline{z} = a - bi$. Đối xứng qua trục $Ox$
\end{dinhnghia}

\begin{tinhchat}
  \begin{itemize}
      \item $\overline{\overline{z}} = z$.
      \item $z + \overline{z} = 2a$.
      \item $\overline{z}\overline{z} = |z|^2 = a^2 + b^2$.
      \item $|\overline{z}| = |z|$
      \item $\dfrac{1}{z} = \dfrac{\overline{z}}{|z|^2}$
      \item $\overline{z_1 + z_2} = \overline{z_1} + \overline{z_2}$
      \item $\overline{z_1 . z_2} = \overline{z_1} . \overline{z_2}$
      \item $\overline{(\dfrac{z_1}{z_2})} = \dfrac{\overline{z_1}}{\overline{z_2}}$
  \end{itemize}
\end{tinhchat}

\subsection{Căn bậc $n$}

\begin{dinhnghia}
  \textbf{Căn bậc $n$} của $z = r(\cos\varphi + i\sin\varphi)$ ($z \ne 0$) có đúng $n$ giá trị:
  \[
    w_k = \sqrt[n]{r}\left(\cos\frac{\varphi + 2k\pi}{n} + i\sin\frac{\varphi + 2k\pi}{n}\right), \quad k = 0, 1, \ldots, n-1.
  \]
\end{dinhnghia}

\begin{vidu}
  Căn bậc hai của $-1$: $w_0 = i$, $w_1 = -i$. Căn bậc ba của 1: $1$, $\omega = e^{2\pi i/3}$, $\omega^2$.
\end{vidu}

\subsection{Đa thức trên trường số phức}

\begin{dinhnghia}
  \textbf{Đa thức} với hệ số trong $\mathbb{C}$: $P(z) = a_n z^n + a_{n-1} z^{n-1} + \cdots + a_1 z + a_0$, $a_i \in \mathbb{C}$, $a_n \ne 0$. Bậc $\deg P = n$.
\end{dinhnghia}

%% ================================================================
%%  CHƯƠNG 3: PHƯƠNG TRÌNH BẬC HAI
%% ================================================================

\section{Phương trình bậc hai}

\begin{dinhnghia}
  Phương trình $az^2 + bz + c = 0$ ($a \ne 0$, $a, b, c \in \mathbb{C}$). Biệt thức $\Delta = b^2 - 4ac$. Trên $\mathbb{C}$, mọi phương trình bậc hai đều có nghiệm (có thể phức).
\end{dinhnghia}

\begin{congthuc}
  Nếu $\Delta \ne 0$: hai nghiệm $z = \dfrac{-b \pm \sqrt{\Delta}}{2a}$. Nếu $\Delta = 0$: nghiệm kép $z = -\dfrac{b}{2a}$.
\end{congthuc}

\begin{dinhly}[title={Định lý cơ bản của đại số}]
  Mọi đa thức bậc $n \ge 1$ với hệ số phức có ít nhất một nghiệm phức. Hệ quả: đa thức bậc $n$ có đúng $n$ nghiệm (kể cả bội) trong $\mathbb{C}$.
\end{dinhly}

\begin{vidu}
  $P(z) = z^2 + 1$ có nghiệm $i$, $-i$. 
  
  $P(z) = z^n - 1$ có $n$ nghiệm là các căn bậc $n$ của đơn vị: $e^{2\pi ik/n}$, $k = 0, \ldots, n-1$.
\end{vidu}

%% ================================================================
%%  CHƯƠNG 4: ĐA THỨC
%% ================================================================

\section{Đa thức}

\begin{tinhchat}[title={Phép toán trên đa thức}]
  Cộng, trừ, nhân đa thức: tương tự số thực. \textbf{Chia đa thức}: $P(x) = Q(x) \cdot D(x) + R(x)$ với $\deg R < \deg Q$. \textbf{Đa thức đồng nhất}: $P(z) = 0$ với mọi $z$ $\Leftrightarrow$ mọi hệ số bằng 0.
\end{tinhchat}

%% ================================================================
%%  CHƯƠNG 5: PHÂN THỨC HỮU TỈ
%% ================================================================

\section{Phân thức hữu tỉ}

\begin{dinhnghia}
  \textbf{Phân thức hữu tỉ}: $\dfrac{P(x)}{Q(x)}$ với $P, Q$ là đa thức, $Q(x) \ne 0$. \textbf{Phân thức thực sự}: $\deg P < \deg Q$. \textbf{Phân thức không thực sự}: $\deg P \ge \deg Q$ (chia được $P = Q \cdot D + R$, đưa về đa thức cộng phân thức thực sự $\dfrac{R}{Q}$).
\end{dinhnghia}

\begin{dinhnghia}
  \textbf{Phân thức thực sự đơn giản}: $\dfrac{A}{(x - a)^k}$ hoặc $\dfrac{Bx + C}{(x^2 + px + q)^m}$ (với $x^2 + px + q$ vô nghiệm).
\end{dinhnghia}

\begin{dinhly}[title={Phân tích thành tổng phân thức đơn giản}]
  Mọi phân thức thực sự $\dfrac{P}{Q}$ (với $Q$ phân tích được) đều viết được thành tổng các phân thức thực sự đơn giản. Phương pháp: đồng nhất hệ số hoặc thay giá trị đặc biệt.
\end{dinhly}

\end{document}
